\documentclass[10pt,a4paper]{amsart}
\usepackage[margin=19mm]{geometry}
\usepackage{amsmath,amssymb,mathtools}
\usepackage[scr=rsfs,cal=euler]{mathalfa}
\usepackage{xfrac}
\usepackage[unicode,hidelinks,bookmarks=false]{hyperref}
\newcommand{\ie}{i.e.\ }
\newcommand{\cf}{cf.\ }
\newcommand{\resp}{respectively\ }
\newcommand{\Subsection}{\subsection}
\providecommand{\nobreakdash}{\nobreak\hbox{-}}
\setlength{\emergencystretch}{2em}
\multlinegap0pt
\usepackage{amssymb}
\usepackage{float}
\newtheorem{thm}{Theorem}
\newtheorem{cor}{Corollary}
\newcommand{\be}{\begin{equation}}
\newcommand{\ee}{\end{equation}}
\title{On Berwald Spaces with non-Zero Flag Curvature}
\author{A. Tayebi, B. Najafi}
\date{Source: arXiv:2601.20087v1 (CC BY 4.0)}
\begin{document}
\maketitle

\noindent\emph{Source and licence.} On Berwald Spaces with non-Zero Flag Curvature. Version arXiv:2601.20087v1 (CC BY 4.0), distributed by arXiv under the Creative Commons Attribution 4.0 International licence, http://creativecommons.org/licenses/by/4.0/, as recorded in the arXiv licence metadata for this exact version and shown on the abstract page https://arxiv.org/abs/2601.20087v1. Full licence notice: \texttt{licenses/arxiv-2601.20087-CC-BY-4.0.txt}.\par\smallskip

\noindent\textit{Editorial scope.} Counted unit: the article's cor CSZ (source0.tex lines 333--335). The article's complete original proof of it is delivered with it (source0.tex lines 336--338, one \texttt{proof} environment of 3 lines). No mathematics is written, repaired or re-derived: the statement and the proof are the article's own lines, in the article's own notation, and no retained line is substituted or re-flowed.  Retained uncounted as the source-local context the proof needs: lines 75--81; lines 242--296; lines 304--326.  Nothing was omitted from any retained span, so the package carries no omission record.  The retained lines cite 4 works, whose entries are transcribed verbatim from the article's own bibliography source in the e-print and delivered with short editorial labels recorded in the item's reference mapping.  No retained line contains a figure environment, an image-inclusion command or drawing code, so no image is delivered and the package declares no asset dependency.  Editorial formatting only, in addition: an AMS \texttt{amsart} preamble that restates the article's own declarations and macros used by the retained lines, namely the environments \texttt{thm}, \texttt{cor} on separate counters, together with the standard packages those lines need.  The retained environments print in this document as Theorem 1, Corollary 1 in order of appearance, whereas the article prints the same items with the numbering of its own full document.  The complete native source of the article as distributed in the arXiv e-print for this exact version is bundled unchanged as \texttt{sources/193504/source0.tex}.
\begin{thm}\label{MTHM}
Let $(M, F)$ be a Berwald manifold. Then, the following hold
\begin{itemize}
\item[1.] If\ \ ${\bf K}(P, y)$ vansihes for all flags and flag-poles, then $F$ is locally Minkowskian;
\item[2.] If\ \ ${\bf K}(P, y)$ is no where zero, then $F$ is Riemannian metric.
\end{itemize}
\end{thm}

%---------------------------------------------------------------------------------------------------------------------------------------------------
\section{Proof of Theorem \ref{MTHM}}
%---------------------------------------------------------------------------------------------------------------------------------------------------
Express the Riemann curvature by
${\bf R}= R^i_{\ k} \omega^k \otimes {\bf e}_i$.
Let
\be
R^i_{\ kl}:= {1\over 3} \Big \{\frac{\partial R^i_{\ k}}{\partial y^l} - \frac{\partial R^i_{\ l}}{\partial y^k} \Big \}.\label{Kikl}
\ee
In fact,  $R^i_{\ kl}$ and $L^i_{\ kl}:= g^{ij}L_{jkl}$ are determined by the following equation
\[
\Omega^i := d\omega^{n+i}-\omega^{n+j} \wedge \omega_j^{\ i}
= {1\over 2} R^i_{\  kl}\omega^k \wedge \omega^l - L^i_{\ kl} \omega^k \wedge \omega^{n+l}.
\]
In this paper, we use the Berwald connection  and the $h$- and $v$- covariant derivatives of a Finsler tensor field are denoted by `` $|$ " and ``, " respectively. In local coordinate system, the Berwald connection determined by following
\begin{eqnarray}
&&d\omega^i = \omega^j \wedge \omega^i_j,\\
&&dg_{ij}- g_{kj}\omega^k_i - g_{ik}\omega^k_j = -2L_{ijk}\omega^k + 2C_{ijk}\omega^{n+k}.\label{Berwaldms}
\end{eqnarray}
Thus
\[
g_{ij|k}= -2L_{ijk},   \ \ \ \ \ g_{ij,k}= 2C_{ijk}.
\]
Also, we have
\begin{equation}
\Omega^i_{\ j}=d\omega^i_{\ j}-\omega^k_{\ j}\wedge\omega^i_{\
k}=\frac{1}{2}R^i_{\ jkl}\omega^k\wedge\omega^l-B^i_{\
jkl}\omega^k\wedge\omega^{n+l}.\label{BCon}
  \end{equation}
Thus
\begin{eqnarray}
&&R^h_{\ ijk}+ R^h_{\ jki}+ R^h_{\ kij}=0,\label{xxxzzz}\\
&&B^h_{\ ijk}=B^h_{\ jki}=B^h_{\ kij}.
\end{eqnarray}
Differentiating (\ref{Berwaldms}) implies that
\begin{eqnarray}
\nonumber g_{pj}\Omega^p_i+g_{ip}\Omega^p_j=\!\!\!\!&-&\!\!\!\!\! 2L_{ijk|s}\omega^k\wedge\omega^s-2L_{ijk,s}\omega^k\wedge\omega^{n+s}\\
\!\!\!\!&-&\!\!\!\!\! 2C_{ijs|k}\omega^k\wedge\omega^{n+s}-2C_{ijs,k}\omega^{n+k}\wedge\omega^{n+s}-2C_{ijs}\Omega^s.\label{mohem2}
\end{eqnarray}
Putting (\ref{BCon}) in (\ref{mohem2}) gives us
\begin{eqnarray}
&&L_{ijk|l}-L_{ijl|k}=-\frac{1}{2}g_{pj}R^p_{i\ kl}-\frac{1}{2}g_{ip}R^p_{j\ kl}-C_{ijl}R^l_{\ kl}.\label{xxxyyy}
\end{eqnarray}
By contracting (\ref{xxxyyy}) with $y^l$, we get
\be
L_{ijk;m}y^m + C_{ijm}R^m_{\ \ k}
= - {1\over 2} g_{im} R^m_{\ \ kl\cdot j} y^l - {1\over 2} g_{jm} R^m_{\ \ kl\cdot j} y^l,\label{LCKK}
\ee
Plugging (\ref{Kikl}) into (\ref{LCKK}) yields
\begin{eqnarray}
L_{ijk|m}y^m + C_{ijm}R^m_{\ \ k}=   - {1\over 3}g_{im}R^m_{\ \ k\cdot j}
- {1\over 3} g_{jm} R^m_{\ \ k\cdot i}- {1\over 6} g_{im}R^m_{\ \ j\cdot k}
- {1\over 6} g_{jm}R^m_{\ \ i\cdot k}. \label{Moeq1}
\end{eqnarray}
For more details about the relation \eqref{Moeq1}, one can refer to \cite{MS}.

Now, by multiplying (\ref{Moeq1}) with $g^{ij}$  we get
\be
J_{k|m}y^m + I_mR^m_{\ k}  = -  {1\over 3} \Big \{ 2 R^m_{\ \; k\cdot m} + R^m_{\ \; m\cdot k} \Big \} \label{Moeq2}
\ee
Also, S-curvature satisfies the following equation
\be
{\bf S}_{\cdot k|m}y^m - {\bf S}_{| k}
= - {1\over 3} \Big \{ 2 R^m_{\ \; k\cdot m} + R^m_{\ \; m\cdot k} \Big \}.\label{SSKKMo}
\ee
For more details, see \cite{CMS}. Comparing (\ref{Moeq2}) and (\ref{SSKKMo}) yields
\be
J_{k|m}y^m +I_m R^m_{\ k} = {\bf S}_{\cdot k|m}y^m - {\bf S}_{|k}. \label{EIILJ}
\ee
See \cite{Shn}. We can rewrite \eqref{EIILJ} as follows
\be
I^i_{\ | p|q} y^py^q +  R^i_{\ m}I^m = g^{ik} \Big \{ {\bf S}_{\cdot k|m}y^m - {\bf S}_{|k} \Big \}. \label{EIILJ*}
\ee
By assumption, ${\bf B}=0$. Every Berwald metric satisfies ${\bf S}=0$ and ${\bf J}=0$. Then, \eqref{EIILJ*} reduces to
\be
{\bf R}_y({\bf I}_y)=0. \label{RI}
\ee
Since ${\bf I}_y$ is orthogonal to $y$ with respect to ${\bf g}_y$, then by \eqref{RI}, it follows that  ${\bf K}(P_0, y)=0$, where
$P_0= {\rm span} \{{\bf I}_y, y\}$ whenever ${\bf I}_y \not=0$. But by assumption, we have ${\bf K}(P_0, y)<0$. This contradiction yields that   ${\bf I}_y =0$ for all $y\in T_xM_0$.

\begin{cor}{\rm (Szab\'{o} Rigidity Theorem)}\label{CSZ}
Let $(M, F)$ be a Berwald surface. Then $F$ is Riemannian metric or locally Minkowskian.
\end{cor}

\begin{proof}
Let $F$ be a Berwald surface (or of scalar flag curvature), $R^i_{\ m}={\bf K}F^2h^i_m$. Then \eqref{RI} reduces to ${\bf K}I_i =0$. Since ${\bf B}=0$, then by Akbar-Zadeh theorem ${\bf K}={\bf K}(x)$ (see \cite{AZ}). Let  ${\bf K}(x)\neq 0$, $\forall x\in M$. Then ${\bf I}=0$ and $F$ is a Riemannian metric.  If $F$ is not Riemannian, then ${\bf K}(x)= 0$. Every Berwald metric with vanishing flag curvature is locally Minkowskian.
\end{proof}

\begin{thebibliography}{99}
\bibitem[AZ]{AZ} H. Akbar-Zadeh, {\it Sur les espaces de Finsler \'{a} courbures sectionnelles constantes}, Bull. Acad. Roy. Bel. Cl, Sci, 5e S\'{e}rie - Tome LXXXIV (1988), 281-322. %--------------------------------------------------------------------------------------------------------------------------------------------------
\bibitem[CMS]{CMS} X. Chen(g), X. Mo and Z. Shen, {\it On the flag curvature of Finsler metrics of scalar curvature}, J. London Math. Soc. {\bf 68}(2003), 762-780. %-------------------------------------------------------------------------------------------------------------------------------------------------- %
\bibitem[MS]{MS} X. Mo and Z. Shen, {\it On negatively curved Finsler manifolds of scalar curvature}, Canad. Math. Bull. {\bf 48}(1) (2005), 112-120. %-------------------------------------------------------------------------------------------------------------------------------------------------
\bibitem[Shn]{Shn} Z. Shen, {\it Nonpositively curved Finsler manifolds with constant S-curvature}, Math. Zeitschrift, {\bf 249}(2005), 625-639. %--------------------------------------------------------------------------------------------------------------------------------------------------
\end{thebibliography}
\end{document}
